\documentclass{article}

\PassOptionsToPackage{numbers, sort, compress}{natbib}
\usepackage[main,final]{neurips_2025}

\usepackage[T2A]{fontenc}
\usepackage[utf8]{inputenc}
\usepackage[english,russian]{babel}
\usepackage{hyperref}
\usepackage{url}
\usepackage{booktabs}
\usepackage{nicefrac}
\usepackage{microtype}
\usepackage{xcolor}
\usepackage{graphicx}
\usepackage{amsmath,amssymb,amsfonts}

\title{Сравнение стратегий нарезки патентных документов\\для семантического поиска и RAG}

\author{
  Елизавета Кожемякова\\
  TODO: организация\\
  \texttt{TODO: email}\\
  \And
  Роман Ищенко\\
  TODO: организация\\
  \texttt{TODO: email}\\
}

\begin{document}

\maketitle

%%%%%%%%%%%%%%%%%%%%%%%%%%%%%%%%%%%%%%%%%%%%%%%%%%%%%%%%%%%%
% Аннотация: актуальность (1) — цель (1) — методы (2–3) — результаты (2) — выводы (1).
% Без ссылок и нерасшифрованных сокращений. ~200 слов.
\begin{abstract}
Качество семантического поиска и генерации с использованием извлеченной информации (Retrieval-Augmented Generation, или RAG) по длинным техническим документам во многом определяется тем, каким образом документ разбивается на фрагменты («чанки») перед векторизацией. Данная работа посвящена поиску схемы нарезки и векторизации патентных документов, которая повышает качество поиска при использовании готовых моделей векторных представлений текста без их дообучения. Патент~--- длинный и неоднородный документ. Его разделы~--- сущность изобретения, уровень техники, подробное описание, формула~--- существенно различаются по информативности, а нарезка фрагментами фиксированной длины смешивает их, размывает смысловой сигнал и позволяет шаблонным формулировкам доминировать в выдаче. Предлагается формировать фрагменты по смысловой структуре самого документа~--- по разделам и по блокам, выделяемым триггерными фразами,~--- а не по длине текста. Структурно-ориентированные схемы сравниваются на корпусе из 37\,727 патентов США с нарезкой фиксированной длины, группировкой абзацев и иерархическим индексированием на основе суммаризации языковой моделью. В качестве экспертной разметки релевантности используются ссылки патентов на документы уровня техники. Кроме того, для каждого раздела оценивается его вклад в качество поиска.
Полученные результаты позволяют определить, какие части патентного документа следует векторизовать в первую очередь. Практическая ценность работы состоит в рекомендациях по выбору схемы нарезки для поиска, кластеризации и RAG по научно-технической информации.
\end{abstract}

\textbf{Ключевые слова:} генерация с использованием извлеченной информации, плотный поиск, нарезка документов, патентные документы, векторные представления текста

%%%%%%%%%%%%%%%%%%%%%%%%%%%%%%%%%%%%%%%%%%%%%%%%%%%%%%%%%%%%
\section{Введение}\label{sec:intro}

TODO

\section{Обзор литературы}\label{sec:rw}

TODO

%%%%%%%%%%%%%%%%%%%%%%%%%%%%%%%%%%%%%%%%%%%%%%%%%%%%%%%%%%%%
\section{Постановка задачи}\label{sec:prelim}

TODO

\section{Метод}\label{sec:method}

TODO

\section{Вычислительный эксперимент}\label{sec:exp}

TODO

\section{Обсуждение}\label{sec:disc}

TODO

\section{Заключение}\label{sec:concl}

TODO

%%%%%%%%%%%%%%%%%%%%%%%%%%%%%%%%%%%%%%%%%%%%%%%%%%%%%%%%%%%%

\bibliographystyle{unsrtnat}
\bibliography{references}

\end{document}
